\documentclass[11pt]{article}
\usepackage[margin=1in]{geometry}
\usepackage{amsmath,amssymb,amsthm}
\usepackage[hidelinks]{hyperref}

\newtheorem{theorem}{Theorem}
\newtheorem{lemma}{Lemma}
\newtheorem{corollary}{Corollary}

\newcommand{\R}{\mathbb R}
\newcommand{\N}{\mathbb N}
\newcommand{\St}{\mathrm{St}}
\newcommand{\Rad}{\mathrm{Rad}}
\newcommand{\cL}{\mathcal L}
\newcommand{\cA}{\mathcal A}
\newcommand{\HT}{\mathcal H}
\newcommand{\norm}[1]{\left\|#1\right\|}

\title{Removing Nontrivial Type from Rozendaal's Endpoint Lattice Case}
\author{Run \texttt{fa\_banach\_001}, agent lane 17}
\date{2026-06-20}

\begin{document}
\maketitle

\section*{Status}

This packet gives a candidate full answer, likely valid, to the question after
Theorem 6.3 in Rozendaal's paper \cite{Rozendaal}.  The source asks whether the
assumption that \(X\) has nontrivial type is necessary in the \(R\)-bounded
operator-valued endpoint calculus.  It is not necessary: the theorem extends to
the endpoint \(p=1\) for complemented subspaces of \(q\)-concave Banach lattices.

\section*{Source Question}

After proving Theorem 6.3, Rozendaal writes:
\[
\text{``We do not know whether the assumption in Theorem 6.3 that \(X\) has
nontrivial type is necessary.''}
\]
The proof there uses nontrivial type only to obtain, uniformly in \(n\), a
projection from \(L^2([0,1];X)\) onto \(\Rad^n(X)\).  We replace that step by a
direct lattice square-function model.

\section*{Idea of the Proof}

Let \(Y\) be a \(q\)-concave Banach lattice, \(2\le q<\infty\).  The finite
Rademacher space \(\Rad^n(Y)\) is uniformly isomorphic to the lattice
\[
        Y(\ell_2^n)
        =\left\{(y_1,\ldots,y_n):
          \norm{(y_k)}_{Y(\ell_2^n)}
          =\norm{\left(\sum_{k=1}^n |y_k|^2\right)^{1/2}}_Y<\infty\right\}.
\]
This is the standard Khintchine square-function equivalence for Banach
lattices of finite cotype.  Moreover \(Y(\ell_2^n)\) is again \(q\)-concave
with constants independent of \(n\), because pointwise
\(\ell_q(\ell_2)\le \ell_2(\ell_q)\) for \(q\ge2\).  Thus if \(X\) is
complemented in \(Y\), then \(\Rad^n(X)\) is uniformly complemented in a
\(1\)-convex, \(q\)-concave Banach lattice.  This is exactly the uniform
geometric input required in Rozendaal's proof of Theorem 6.3.

\section*{The Uniform Rademacher Lattice Model}

\begin{lemma}\label{lem:rad-lattice}
Let \(Y\) be a \(q\)-concave Banach lattice for some \(2\le q<\infty\).  Then
for every \(n\in\N\), \(\Rad^n(Y)\) is isomorphic, with constants independent
of \(n\), to a \(1\)-convex and \(q\)-concave Banach lattice.  More precisely,
\(\Rad^n(Y)\) is uniformly isomorphic to \(Y(\ell_2^n)\), and
\(Y(\ell_2^n)\) has \(q\)-concavity constant at most the \(q\)-concavity
constant of \(Y\).
\end{lemma}

\begin{proof}
Since \(Y\) has finite cotype, the standard Khintchine inequality for Banach
lattices gives constants independent of \(n\) such that
\[
        \left(\mathbb E\norm{\sum_{k=1}^n r_k y_k}_Y^2\right)^{1/2}
        \simeq
        \norm{\left(\sum_{k=1}^n |y_k|^2\right)^{1/2}}_Y
\]
for all \(y_1,\ldots,y_n\in Y\); see, for instance, the lattice cotype
material in \cite{LT}.  Hence \(\Rad^n(Y)\) is uniformly isomorphic to
\(Y(\ell_2^n)\).

It remains to record that the lattice constants do not depend on \(n\).  Every
Banach lattice is \(1\)-convex.  Let \(v^j=(y_{j,k})_{k=1}^n\in Y(\ell_2^n)\),
\(j=1,\ldots,m\), and set
\[
        s_j=\left(\sum_{k=1}^n |y_{j,k}|^2\right)^{1/2}\in Y_+.
\]
If \(C_q(Y)\) is the \(q\)-concavity constant of \(Y\), then
\[
        \left(\sum_{j=1}^m \norm{v^j}_{Y(\ell_2^n)}^q\right)^{1/q}
        =
        \left(\sum_{j=1}^m \norm{s_j}_{Y}^q\right)^{1/q}
        \le
        C_q(Y)\norm{\left(\sum_{j=1}^m s_j^q\right)^{1/q}}_Y.
\]
Pointwise, for nonnegative scalars \(a_{j,k}\),
\[
        \left(\sum_j\left(\sum_k a_{j,k}^2\right)^{q/2}\right)^{1/q}
        \le
        \left(\sum_k\left(\sum_j a_{j,k}^q\right)^{2/q}\right)^{1/2},
\]
which is Minkowski's inequality in the form
\(\ell_q(\ell_2)\hookrightarrow \ell_2(\ell_q)\) for \(q\ge2\).  Applying this
inside the lattice \(Y\) gives
\[
        \norm{\left(\sum_j s_j^q\right)^{1/q}}_Y
        \le
        \norm{\left(\sum_{k=1}^n
        \left(\sum_{j=1}^m |y_{j,k}|^q\right)^{2/q}\right)^{1/2}}_Y
        =
        \norm{\left(\sum_{j=1}^m |v^j|^q\right)^{1/q}}_{Y(\ell_2^n)}.
\]
Thus \(Y(\ell_2^n)\) is \(q\)-concave with constant at most \(C_q(Y)\),
uniformly in \(n\).
\end{proof}

\begin{lemma}\label{lem:complemented-rad}
Let \(X\) be isomorphic to a complemented subspace of a \(q\)-concave Banach
lattice \(Y\), where \(2\le q<\infty\).  Then \(\Rad^n(X)\), uniformly in
\(n\), is isomorphic to a complemented subspace of a \(1\)-convex,
\(q\)-concave Banach lattice.
\end{lemma}

\begin{proof}
Identify \(X\) with a complemented subspace of \(Y\), and let \(P:Y\to X\) be a
bounded projection.  The coordinatewise map
\[
        \widetilde P\left(\sum_{k=1}^n r_k y_k\right)
        =
        \sum_{k=1}^n r_k P y_k
\]
is a projection from \(\Rad^n(Y)\) onto \(\Rad^n(X)\), with
\(\norm{\widetilde P}\le\norm{P}\).  Now apply Lemma \ref{lem:rad-lattice}.
\end{proof}

\section*{Endpoint \(R\)-Bounded Calculus}

\begin{theorem}\label{thm:endpoint}
Let \(X\) be isomorphic to a complemented subspace of a \(q\)-concave Banach
lattice for some \(2\le q<\infty\).  Let \(-iA\) generate a \(C_0\)-group
\((U(s))_{s\in\R}\subseteq\cL(X)\), let \(\lambda>\omega>\theta(U)\), and let
\(\cA\subseteq\cL(X)\) be the algebra of bounded operators commuting with \(A\).
Then there exists a constant \(C\ge0\) such that
\[
 R_{D((\lambda+iA)^{1-1/q}),X}
 \left(\{f(A): f\in R\HT^\infty(\St_\omega;\cA\cap\mathcal T)\}\right)
 \le C\,R_X(\mathcal T)
\]
for every \(R\)-bounded set \(\mathcal T\subseteq\cL(X)\).
\end{theorem}

\begin{proof}
We follow Rozendaal's proof of Theorem 6.3 \cite{Rozendaal}, changing only the
uniform geometric input.  Put \(\theta=1-1/q\).  Fix \(n\in\N\), let
\(\Rad^n(X)\subset L^2([0,1];X)\), and let \(\widetilde A\) be the diagonal
operator
\[
        \widetilde A\left(\sum_{k=1}^n r_k x_k\right)
        =\sum_{k=1}^n r_k A x_k .
\]
Then \(-i\widetilde A\) generates the diagonal group on \(\Rad^n(X)\).

By Lemma \ref{lem:complemented-rad}, \(\Rad^n(X)\) is uniformly isomorphic to a
complemented subspace of a \(1\)-convex, \(q\)-concave Banach lattice.  Hence
the operator-valued fractional-domain estimate preceding Theorem 6.2 in
\cite{Rozendaal}, applied to \(\widetilde A\), gives a constant independent of
\(n\) such that
\[
        \norm{F(\widetilde A)}_
        {\cL(D((\lambda+i\widetilde A)^\theta),\Rad^n(X))}
        \le C_0\norm{F}_{R\HT^\infty(\St_\omega;\widetilde{\cA})}
\]
for every admissible operator-valued \(F\).  The independence of \(n\) follows
from Rozendaal's Remark 5.3, because all complementability and \(q\)-concavity
constants supplied by Lemma \ref{lem:complemented-rad} are uniform in \(n\).

Now take \(f_1,\ldots,f_n\in R\HT^\infty(\St_\omega;\cA\cap\mathcal T)\) and
define
\[
        F(z)\left(\sum_{k=1}^n r_k x_k\right)
        =
        \sum_{k=1}^n r_k f_k(z)x_k .
\]
Since \(X\) is a complemented subspace of a finite-cotype Banach lattice, \(X\)
has property \((\alpha)\), as recalled in \cite{Rozendaal}.  The same
property-\((\alpha)\) estimate used in the proof of Rozendaal's Theorem 6.2
shows
\[
        \norm{F}_{R\HT^\infty(\St_\omega;\widetilde{\cA})}
        \le C_1 R_X(\mathcal T),
\]
with \(C_1\) independent of \(n\), \(f_k\), and \(\mathcal T\).

Finally, by the diagonal functional calculus,
\[
        F(\widetilde A)\left(\sum_{k=1}^n r_k x_k\right)
        =
        \sum_{k=1}^n r_k f_k(A)x_k,
\]
and
\[
        D((\lambda+i\widetilde A)^\theta)
        =
        \Rad^n\!\left(D((\lambda+iA)^\theta)\right)
\]
with the natural graph norms.  Therefore
\[
        \left\|\sum_{k=1}^n r_k f_k(A)x_k\right\|_{L^2([0,1];X)}
        \le
        C_0C_1 R_X(\mathcal T)
        \left\|\sum_{k=1}^n r_k x_k\right\|_
        {L^2([0,1];D((\lambda+iA)^\theta))}.
\]
This is exactly the asserted \(R\)-bound.
\end{proof}

\begin{corollary}
The nontrivial type assumption in Rozendaal's Theorem 6.3 is not necessary.
\end{corollary}

\begin{proof}
Take, for example, \(X=L^1[0,1]\).  This Banach lattice is \(q\)-concave for
each \(q\ge2\), but it has no nontrivial type.  Theorem
\ref{thm:endpoint} applies to \(X\), giving the endpoint \(R\)-bounded calculus
with exponent \(1-1/q\).  Thus the conclusion of the theorem can hold without
nontrivial type.
\end{proof}

\section*{Verification Notes}

The proof uses only the source theorem and standard Banach-lattice facts.  The
new step is Lemma \ref{lem:rad-lattice}; its constants are independent of
\(n\), which is precisely what the original proof needed from \(K\)-convexity.
The argument does not claim that every Banach space with property
\((\alpha)\) and finite cotype admits the endpoint theorem.  It answers the
source question by proving that nontrivial type is not necessary: finite-cotype
lattices such as \(L^1\) already provide no-nontrivial-type cases.

\section*{Novelty Check}

The run indexes were searched for \texttt{1508.02036}, \texttt{Functional
calculus for C\_0-groups}, \texttt{second R-bounded calculus},
\texttt{nontrivial type}, \texttt{K-convex}, \texttt{Rademacher projection},
and related transference/Fourier multiplier phrases.  Web searches on
2026-06-20 for exact phrases from the source question and theorem found the
source paper and related Rozendaal--Veraar multiplier papers, but no later
answer to this endpoint necessity question.

\begin{thebibliography}{9}

\bibitem{Rozendaal}
J. Rozendaal,
\emph{Functional calculus for \(C_0\)-groups using (co)type},
arXiv:1508.02036, 2015.

\bibitem{LT}
J. Lindenstrauss and L. Tzafriri,
\emph{Classical Banach Spaces II: Function Spaces},
Ergebnisse der Mathematik und ihrer Grenzgebiete 97,
Springer-Verlag, 1979.

\end{thebibliography}

\end{document}
